\documentclass[12pt]{article}
\usepackage[utf8]{inputenc}
\usepackage[T1]{fontenc}
\usepackage{hyperref}
\title{Iyansi Language Resources for K-12 Curriculum Development}
\author{Compiled by Claude Sonnet 5.5 (claude-sonnet-5-5) for Andr\'as Kornai}
\date{}
\begin{document}
\maketitle

\section{Introduction}
Iyansi (Yansi, Eyansi, Kiyansi, Yanzi; ISO 639-3: yns; Glottocode yans1239) is a Bantu language of southwestern DR Congo, around Bandundu (Guthrie B.85). Speakers call themselves Bayanzi. The figure of about 100,000 speakers dates from 1997. Wikitongues gives its EGIDS status as 6a (vigorous). Listed varieties include Mbiem (West Yansi), East Yans, Yeei, Tsong, Mpur and Tsambaan. Sources disagree on classification and on whether Tsong/Songo is a separate language.

\section{Primary Educational Resources}
No primers, textbooks or school programmes were found in the searches.

\section{Linguistic and Reference Materials}
\begin{itemize}
\item Wikipedia, ``Yansi language'': classification, varieties and orthography notes. \url{https://en.wikipedia.org/wiki/Yansi_language}
\item Omniglot, ``Yansi (Eyansi)'': notes on the writing system. A Latin-script orthography dates from 1985. \url{https://www.omniglot.com/writing/yansi.htm}
\item Wikitongues language page. \url{https://wikitongues.org/languages/yns/}
\end{itemize}
No grammar or dictionary was found in the search results.

\section{Religious and Community Resources}
\begin{itemize}
\item Joshua Project, Yansi in DR Congo: people-group profile. \url{https://joshuaproject.net/people_groups/15986/CG}
\item Wikipedia, ``Yanzi people''. \url{https://en.wikipedia.org/wiki/Yanzi_people}
\end{itemize}
Wikipedia reports a New Testament translation exists in Yansi, but the search results did not show a link to the text itself.

\section{Digital and Audio Resources}
\begin{itemize}
\item Global Recordings Network, Yansi (yns) recordings page. \url{https://globalrecordings.net/en/language/yns}
\end{itemize}

\section{Curriculum Development Context}
Iyansi is vigorous, with a 1985 Latin orthography and a New Testament, but no school materials were found. A curriculum would build on the New Testament and orthography and would need new materials and community partners. Variation among the dialects should be addressed early. No contact details were found, so none are given.

\end{document}